\documentclass{exam}
\usepackage{amsmath}
\usepackage{graphicx}
\usepackage{nicematrix}
\usepackage{extpfeil}
\usepackage{amsthm}
\usepackage{gensymb}
\usepackage{amsfonts}
\usepackage{tikz}
\usepackage{pgfplots}
\usepackage{array}
\usepackage{hyperref}
\usepackage{nameref}
\usepackage[x11names]{xcolor}
\usepackage[most]{tcolorbox}

\pgfplotsset{compat=1.18}
\printanswers

\newcommand{\ZZ}{\mathbb Z}
\newcommand{\QQ}{\mathbb Q}
\newcommand{\NN}{\mathbb N}
\newcommand{\CC}{\mathbb C}
\newcommand{\RR}{\mathbb R}
\newcommand{\braces}[1]{\ensuremath{\left\{#1 \right\}}}
\newcommand{\Span}[1]{\ensuremath\text{Span}\braces{#1}}
\newcommand{\ee}{\mathbf{e}}

\hypersetup{colorlinks=true, linktoc=section, linkcolor=blue}

\pagestyle{headandfoot}
\firstpageheadrule
\runningheadrule
\firstpageheader{Prof. Shen \\ Differential Equations}{Challenge Problems \#6}{Jeevan Shah}
\runningheader{Differential Equations \\ Challenge Problems \#6}{}{Shah}
\firstpagefooter{}{\thepage}{}
\runningfooter{ }{\thepage}{ }

\printanswers


\begin{document}
    \begin{questions}
        \question Show that if $A$ is self-adjoint, then $e^{tA}$ is contractive if all eigenvalues of $A$ are nonpositive. 
        \begin{solution}
            \begin{proof}
                Assume $A$ is self-adjoint with nonpositive eigenvalues. Since $A$ is self-adjoing, $A^{T} = \bar{A}$, so by the spectral theorem, $A$ is diagonalizable by a amtrix consisting of an orthonormal basis. For $A, n\times\,n$, let $\lambda_1, \ldots, \lambda_n$ be the eigenvalues of $A$ with $\vec{u}_1, \ldots, \vec{u}_n$ the corresponding eigenvalues. Recall that $\lambda_i \leq 0$ for all $i = 1, \ldots, n$ and $\vec{u}_i \cdot \vec{u}_j = 0$ and $\|\vec{u}_i\| = 1$ for $i, j = 1, \ldots, n (i \neq j)$. Then, 
                \[
                    A = PDP^{-1}, \quad P = \begin{pmatrix}
                        \vec{u}_1 & \cdots & \vec{u}_n
                    \end{pmatrix}, \quad D = \begin{pmatrix}
                        \lambda_1 & \cdots & 0 \\
                        \vdots & \ddots & \vdots \\
                        0 & \cdots & \lambda_n 
                    \end{pmatrix}.
                \]
                Since $P$ forms a basis for $\CC^n$ and $\vec{x} \in \CC^n$ we have 
                \[
                    \vec{x} = c_1\vec{u_1} + \cdots + c_n\vec{u_n}
                \]
                for some $c_1, \ldots, c_n \in \CC$. Then, 
                \[
                    e^{tA}\vec{x} = e^{t\lambda_1}c_1\vec{u_1} + \cdots + e^{t\lambda_n}c_n\vec{u_n}
                \]
                so
                \[
                    \|\vec{x}\|^2 = c_1^2 + \cdots + c_n^2 \quad\text{and}\quad \|e^{tA}\|^2 = e^{t\lambda_1}c_1^2 + \cdots + e^{t\lambda_n}c_n^2.
                \]
                Since $e^{t\lambda_i}c_{i}^2 \leq c_i^2$ because $\lambda_i \leq 0$ for $i = 1, \ldots, n$, we have 
                \[
                    \|\vec{x}\| \leq \|e^{tA}\vec{x}\|
                \]
                for all $t$.
            \end{proof}
        \end{solution}

        \question Show that if $K$ is a skew-adjoint square matrix, then $e^{tK}$ is norm-preserving, that is, $\|e^{tK}x\| = \|x\|$ for all $x$ and all $t \geq 0$. Furthermore, show that $e^{tK}$ is unitary, that is, $\left(e^{tK}\right)^{-1} = \left(e^{tK}\right)^{*}$.
        \begin{solution}
            \begin{proof}
                Let $K$ be a skew-adjoint matrix so $K^{*} = -K$, by definition. We first show that $e^{tK}$ is unitary and then that is it norm-preserving. Since $K^{*} = -K$, we have 
                \[
                    \left(e^{tK}\right)^{*} = e^{tK^{*}} = e^{-tK}. 
                \]
                ALos, $K$ communites with $-K$, so 
                \[
                    e^{tK}e^{-tK} = e^{tK - tK} = e^{0} = I. 
                \]
                Thus, $e^{-tK}$ is the inverse of $e^{tK}$. It follows that 
                \[
                    \left(e^{tK}\right)^{*} = \left(e^{tK}\right)^{-1}
                \]
                so $e^{tK}$ is unitary. Now, let $x \in \CC^n$. Since $e^{tK}$ is unitary, we have that 
                \begin{align*}
                    & \| e^{tK}x \|^2 = \langle e^{tK}x, e^{tK}x \rangle = \langle x, \left(e^{tK}\right)^{*}e^{tK}x\rangle = \langle x, Ix \rangle = \langle x, \rangle x = \| x \|^2 \\
                    \Rightarrow& \| e^{tK}x \| = \| x \|
                \end{align*}
                thus, $e^{tK}$ is norm-preserving for all $x \in \CC^n$ and $t \geq 0$.
            \end{proof}
        \end{solution}

        \question Applying the decomposition $A = H + K$, use the Trotter product formula to show that $e^{tA}$ is contractive if $H$ has only non-postive real eigenvalues. 
        \begin{solution}
            \begin{proof}
                Let $A = H + K$ for $H = \frac{1}{2}\left(A + A^*\right)$ and $K = \frac{1}{2}\left(A - A^*\right)$ such that $H$ is self-adjoint and $K$ is skew-adjoint. Assume as well that $H$ has only non-postive eigenvalues. We will show that $e^{tA}$ is contractive for every $t \geq 0$, or equivalently, $\| e^{tA} \| \leq 1$. Since $H$ is self-adjoint and all of its eigenvalues are nonpositive, it follows from exercise $1$ that $e^{tH}$ is contractive so $\| e^{tH} \| \leq 1$ for all $t \geq 0$. By exercise $2$, we have that $e^{tK}$ is norm-preserving and thus unitary since $K$ is skew-adjoint. It follows that $\|e^{tK}\| = 1$ for $t \geq 0$. Now, fix some $t \geq 0$ so by the Trotter Product , 
                \[
                    e^{t(H+K)} = \lim_{n\to\infty}\left(e^{\frac{t}{n}H}e^{\frac{t}{n}K}\right)^n \Rightarrow \| e^{t(H+K)} \| = \lim_{n\to\infty} \|(e^{\frac{t}{n}H}e^{\frac{t}{n}K})^{n}\|
                \]
                by taking operator norms and its continuity over limits. Then, 
                \[
                    \|\left(e^{\frac{t}{n}H} e^{\frac{t}{n}K}\right)^{n}\| \leq \| e^{\frac{t}{n}H} e^{\frac{t}{n}K} \|^{n} \leq \left(\|e^{\frac{t}{n}H}\| \|e^{\frac{t}{n}K}\|\right)^{n}.
                \]
                Since $\|e^{\frac{t}{n}H} \leq 1$ and $\|e^{\frac{t}{n}K} = 1$, we have that 
                \[
                    \left(\|e^{\frac{t}{n}H}\| \|e^{\frac{t}{n}K}\|\right)^{n} \leq 1^{n} = 1
                \]
                so $\| (e^{\frac{t}{n}H} e^{\frac{t}{n}K})^n \| \leq 1$. Passing the limit gives 
                \[
                    \|e^{t(H+K)}\| \leq 1 \Rightarrow \| e^{tA} \| \leq 1
                \]
                so $e^{tA}$ is contractive. 
            \end{proof}
        \end{solution}

        \question Suppose that $A$ is any constant square matrix and that $e^{tA}$ is contractive for all $t \geq 0$. Show then that $A$ is dissipative. 
        \begin{solution}
            \begin{proof}
                Assume that $A$ is a constant square matrix and that $e^{tA}$ is contractive for all $t \geq 0$. We must show that $A$ is dissipative, or equivalently
                \[
                    \Re\langle x, Ax \rangle \leq 0 
                \]
                for all $x \in \CC^n$. We start by fixing $x \in \CC^n$ so that 
                \begin{align*}
                    \|e^{tA}x\| \leq \|x\| &\Rightarrow \| e^{tA}\| ^2 \leq \|x\|^2 \Leftrightarrow \langle e^{tA}x, e^{tA}x \rangle \leq \langle x, x \rangle 
                \end{align*}
                since $e^{tA}$ is contractive. Now, subtractive $\langle x, x \rangle$ from both sides and dividing by $t$ we see that 
                \[
                    \frac{\|e^{tA}x\|^2 - \|x\|^2}{t} \leq 0.
                \]
                Consider, 
                \begin{align*}
                    &\|e^{tA}x\|^2 - \|x\|^2 = \langle e^{tA}x - x, e^{tA}x\rangle + \langle x, e^{tA}x - x\rangle \\
                    \Rightarrow&\, \frac{\|e^{tA}x\|^2 - \|x\|^2}{t} = \left\langle \frac{e^{tA} - I}{t}x, e^{tA}x\right\rangle + \left\langle x, \frac{e^{tA} - I}{t}x\right\rangle.
                \end{align*}
                Now, taking the limit as $t \to 0^{+}$ and applying the facts that 
                \[
                    \lim_{t \to 0^+} e^{tA}x = x \quad\text{and}\quad \lim_{t \to 0^+} \frac{e^{tA} - I}{t}x = Ax 
                \]
                we have that 
                \[
                    \lim_{t \to 0^+} \frac{\|e^{tA}x\|^2 - \|x\|^2}{t} = \langle Ax, x \rangle + \langle x, Ax \rangle.
                \]
                Since $\langle Ax, x \rangle = \overline{\langle x, Ax \rangle}$, it follows that 
                \[
                    \langle Ax, x \rangle + \langle x, Ax \rangle = 2\Re\langle x, Ax \rangle.
                \]
                On the other hand, because 
                \[
                    \frac{\|e^{tA}x\|^2 - \|x\|}{t} \leq 0  
                \]
                for all $t > 0$, passing the limit gives $2\Re\langle x, Ax\rangle \leq 0$, hence, $\Re\langle x, Ax\rangle \leq 0$ for all $x \in \CC^n$. Therefore, $A$ is dissipative.
            \end{proof}
        \end{solution}

        \question Consider 
        \[
            A_1 = \begin{pmatrix}
                -1/3 & 1 \\
                0 & -1/3
            \end{pmatrix} 
            \quad\text{and}\quad A_2 = \begin{pmatrix}
                -1 & 1 \\
                0 & -1
            \end{pmatrix}.
        \]
        Show that althought $A_1$ has only negative eigenvalues, $A_1$ is not dissipative, and so $e^{tA_{1}}$ is not contractive. Show that $A_2$ is dissipative, and so $e^{tA_{2}}$ is contractive. 
        \begin{solution}
            \begin{proof}
            Consider 
            \[
                A_1 = \begin{pmatrix}
                    -1/3 & 1 \\
                    0 & -1/3
                \end{pmatrix} 
                \quad\text{and}\quad A_2 = \begin{pmatrix}
                    -1 & 1 \\
                    0 & -1
                \end{pmatrix}.
            \]
            We will show that $A_1$ has only negative eigenvalues but is not dissipative while $A_2$ is. Since $A_1$ is an upper triangular matrix, we know the eigenvalues of $A_1$ are $\lambda_1 = \lambda_2 = -1/3$, and so, $A_1$ has only negative eigenvalues. Now let 
            \[
                H_1 = \frac{1}{2}\left(A_1 + A_1^*\right) = \frac{1}{2}\begin{pmatrix}
                    -1/3 & 1 \\
                    0 & -1/3
                \end{pmatrix}
                + \frac{1}{2}\begin{pmatrix}
                    -1/3 & 0 \\
                    1 & -1/3
                \end{pmatrix}
                = \begin{pmatrix}
                    -1/3 & 1/2 \\
                    1/2 & -1/3    
                \end{pmatrix}
            \]
            such that $H_1$ is self-adjoint by definition. Then, 
            \[
                \det H_1 = \left(\frac{-1}{3}\right)^2 - \left(\frac{1}{2}\right)^2 = \frac{1}{9} - \frac{1}{4} < 0
            \]
            so the eigenvalues of $H_1$ have opposite signs, or equivalently, $H_1$ has one postive eigenvalue. It follows that there exists $x \in \CC^2$ such that $\langle x, H_1 x\rangle > 0$, but 
            \[
                \Re\langle x, Ax\rangle = \left\langle x, \frac{1}{2}\left(A_1 + A_1^*\right)x \right\rangle = \langle x, Hx \rangle > 0.
            \]
            From exercise $4$, if $e^{tA_{1}}$ were contractive for all $t \geq 0$, then $A_1$ would have to be dissipative. Since $A_1$ is not dissipative, it follows that $e^{tA_{1}}$ is not contractive. Now consider $A_2$ and define 
            \[
                H_2 = \frac{1}{2}\left(A_2 + A_2^*\right) = \begin{pmatrix}
                    -1 & 1/2 \\
                    1/2 & -1 
                \end{pmatrix}
            \]
            such that $H_2$ is also self-adjoint by definition. The eigenvalues of $H_2$ are given by 
            \[
                \lambda^2 + 2\lambda + \frac{3}{4} = \left(\lambda + \frac{1}{2}\right)\left(\lambda + \frac{3}{4}\right) = 0 \Rightarrow \lambda_1 = -\frac{1}{2}, \lambda_2 = -\frac{3}{2}.
            \]
            Since $\lambda_1, \lambda_2 \leq 0$, $H_2$ is self-adjoint with non-postive eignevalues so 
            \[
                \langle x, H_2 x\rangle \leq 0 \Rightarrow \Re\langle x, A_{2}x \rangle = \langle x, H_{2}x \rangle \leq 0 
            \]
            for all $x \in \CC^2$. It follows that $A_2$ is dissipative. Then, by exercise $3$, $e^{tA_{2}}$ is contractive for all $t \geq 0$ since $H_2$ has eigenvalues all nonpositive.
            \end{proof}
        \end{solution}
    \end{questions}

\footnotetext{\LaTeX\, code for this document can be found on github \href{https://github.com/jeevanshah07/MATH292/blob/main/challengeProblems/challengeProblems6/main.tex}{\underline{here}}}
\end{document}